% English translation of capitulos/06-conclusao.tex, made by Claude Code (Claude Fable 5.2) on 24 November 2026.
% Pending review by the author. When the Portuguese file changes, update this one (scripts/verificar_traducao.py).
\section{Conclusion}
This work organized the evidence on prompting techniques for code generation, translation, refactoring and
debugging. To do so, it carried out a rapid literature review that started from 2,304 unique studies in three databases and reached a
corpus of 332 studies included by the author at full text, all with an extraction record. The results
answer the three research questions with a caveat that holds for the whole text: the extraction was done by AI,
and the author
checked 67 studies selected for the sample, in which 18 of 571 blocks required correction, and the 120 empirical claims against the original, of which 30 called for correction.

As for RQ1, the techniques evaluated combine mechanisms, and most of the evidence compares the complete technique with
a baseline. In the extracted comparisons, many favourable results involve verifiable external
information: tests,
execution, repository context and localization of the defect. These gains depend on conditions that are not always
repeated outside the study, such as available tests, more calls to the model or information taken from the reference
solution. Reasoning instructions, self-critique without execution, fixed examples and changes in wording had
mixed or unfavourable results, especially with small models and in repair.

As for RQ2, functional correctness is the most recorded risk, and passing tests is not always equivalent to correct
code. Almost all studies propose some mitigation, but fewer than half evaluate at least one. Contamination and
non-determinism are treated as design care, almost without evaluation, and security instructions may reduce
correctness.

As for RQ3, the multi-label taxonomy separates task, mechanism, role, type of counterfactual and conditions. The decision-support
table summarizes, per task and grouping, the direction of the comparisons, the strength of the evidence and the caveats. It also
marks the cases in which the evidence is not enough to recommend a technique.

In the examination in eight scenarios, the author judged the table partial on most criteria. It led to an executable
script only when the condition of the scenario coincided with that of the studies (D1), and did not declare the evidence
insufficient when reproduction of the defect or tests were missing (D2 and T2). For this reason, the author decided to revise it. Version 2
separates the types of verification, states the result measured in each row and declares in rows of their own the
conditions without sufficient evidence. The inclusion of stage 2 remains to be defined with the advisor.
Outside the sample, the records remain as extraction by AI. The
findings hold for this corpus, which contains only studies with the task in the title and an indication of peer review, and for
results on benchmarks. They do not demonstrate gains in productivity, security or cost in real use.
